\section{РЕЗУЛЬТАТЫ ЭКСПЕРИМЕНТОВ И АНАЛИЗ}

В разделе приведены результаты двух серий экспериментов --- пакетной аналитики на
наборе TPC-DS и near-real-time-сценария при потоковой записи, --- их анализ по типам
запросов и рекомендации по применению форматов хранения.

\subsection{Сценарий 1: пакетная аналитика на наборе TPC-DS}

На наборе TPC-DS~\cite{tpcds-spec} масштаба \SI{100}{\giga\byte} сравнивались
четыре формата --- Parquet, ORC, Vortex и Lance --- по времени выполнения запросов
(таблицы только для добавления, крупные файлы разделов, параллелизм 300); приведены
средние значения по трём прогонам. На всех форматах успешно выполнились
\num{104} запроса. Суммарное время и число «побед» (формат показал наименьшее время)
приведены в таблице~\ref{tab:tpcds-total}.

\begin{table}
    \caption{Суммарное время выполнения запросов TPC-DS (среднее по трём прогонам)}
    \label{tab:tpcds-total}
    \begin{tabularx}{\textwidth}{|X|c|c|}\hline
    \textbf{Формат} & \textbf{Суммарное время, с} & \textbf{Запросов с наименьшим временем} \\ \hline
    Vortex  & 1714 & 90 \\ \hline
    Parquet & 2290 & 5 \\ \hline
    ORC     & 2508 & 1 \\ \hline
    Lance   & 3341 & 8 \\ \hline
    \end{tabularx}
\end{table}

Vortex --- быстрейший формат: по суммарному времени он опережает Parquet на
\SI{25}{\percent}, ORC --- на \SI{32}{\percent}, и показывает наименьшее время на
\num{90} запросах из \num{104}. По числу запросов Vortex быстрее Parquet более чем
на \SI{10}{\percent} на \num{85} запросах, быстрее ORC --- на \num{91}; ORC
опережает Vortex более чем на \SI{10}{\percent} лишь на одном запросе. На восьми
запросах быстрейшим оказался Lance --- несмотря на наихудшее суммарное время: это
запросы инвентаризации с многократным сканированием (q21, q22, q39), где сказывается
его оптимизация под произвольный доступ; на остальных запросах Lance существенно
медленнее. Примечательно, что на пакетной аналитике ORC оказался медленнее Parquet
(Parquet быстрее на \num{74} запросах из \num{104}): преимущество ORC, проявлявшееся
в потоковом сценарии на точечных и узкодиапазонных запросах, не переносится на
полные сканирования крупных файлов, где зрелый векторизованный читатель Parquet
эффективнее. Характерные запросы и ускорение Vortex относительно Parquet и ORC
приведены в таблице~\ref{tab:tpcds}.

\begin{xltabular}{\textwidth}{|l|c|c|X|}
    \caption{Характерные результаты на наборе TPC-DS (масштаб \SI{100}{\giga\byte}): ускорение Vortex относительно Parquet и ORC (среднее по трём прогонам)}\label{tab:tpcds} \\ \hline
    \textbf{Запрос} & \textbf{к Parquet} & \textbf{к ORC} & \textbf{Особенности запроса (объяснение)} \\ \hline
    \endfirsthead
    \caption*{Продолжение таблицы~\ref{tab:tpcds}} \\ \hline
    \textbf{Запрос} & \textbf{к Parquet} & \textbf{к ORC} & \textbf{Особенности запроса (объяснение)} \\ \hline
    \endhead
    \hline
    \endfoot
    \hline
    \endlastfoot
    q51 & $-83\,\%$ & $-83\,\%$ & оконные нарастающие суммы по продажам, мало столбцов, селективные условия --- column pruning и нативное проталкивание предикатов работают в полную силу \\ \hline
    q35 & $-50\,\%$ & $-53\,\%$ & демографический срез с подзапросами существования --- многократное сканирование немногих столбцов \\ \hline
    q88 & $-36\,\%$ & $-45\,\%$ & восемь подсчётов с разными временными окнами --- восьмикратное сканирование одного столбца, выигрыш от column pruning умножается \\ \hline
    q42 & $-36\,\%$ & $-38\,\%$ & агрегация продаж по категории с фильтром по дате --- селективный предикат + узкая проекция \\ \hline
    q73 & $-37\,\%$ & $-38\,\%$ & подсчёт покупателей по числу покупок с фильтрами --- предикат отсекает данные до агрегации \\ \hline
    q1, q2, q28, q83, q99 & медленнее & медленнее & агрегация без селективного предиката над широкими наборами столбцов --- работа в движке агрегации (JVM), а не в декодере \\ \hline
\end{xltabular}

Условия сценария максимально благоприятны для Vortex: таблицы
только для добавления (нет слияния версий), крупные файлы (накладные расходы границы
JNI амортизируются на больших пакетах), реальные данные с избыточностью (кодеки
FastLanes, ALP, FSST дают высокую степень сжатия), запросы с селективными
предикатами и узкими проекциями (нативное проталкивание предикатов сокращает объём
данных, пересекающих границу JVM/native). Пиковое ускорение в \SI{83}{\percent} на
q51 и линейный рост выигрыша с числом сканирований немногих столбцов (q88 --- восемь
сканирований одного столбца) прямо подтверждают первую гипотезу работы.

Запросы, на которых быстрее Parquet, объединяет общая черта --- преобладание
агрегации без селективного предиката над относительно широкими наборами столбцов:
основная работа выполняется не в декодере, а в движке агрегации, который для Parquet
полностью находится в JVM, а для Vortex отделён границей JNI. Это согласуется с
третьей гипотезой и с обсуждавшимся в подразделе~\ref{subsec:methodology}
архитектурным ограничением --- отсутствием проталкивания агрегатов в интерфейсе
форматов Apache Paimon.

Помимо времени запросов, форматы различаются по занимаемому объёму хранилища.
Размер таблиц после полной загрузки приведён в таблице~\ref{tab:tpcds-size}. ORC даёт самый компактный результат (за счёт зрелых
тонких настроек сжатия и эффективного словарного кодирования на категориальных
столбцах TPC-DS); Parquet и Vortex эквивалентны по объёму. Lance занимает в 7--8
раз больше остальных форматов --- следствие оптимизации под произвольный доступ к
строкам (фиксированные страницы, менее агрессивное сжатие), что лишний раз
подтверждает его неприменимость к OLAP-нагрузке. Любопытная особенность Vortex ---
в два раза большее число объектов в хранилище (1206 против 618 у Parquet/ORC) при
том же суммарном объёме: на каждый файл данных приходится дополнительный объект
(сайдкар с переносимыми метаданными). На стоимость хранения это не влияет, но
удваивает число операций \texttt{PUT} в объектное хранилище при записи и число
записей о файлах в манифесте Paimon.

\begin{table}
    \caption{Объём хранилища таблиц TPC-DS масштаба \SI{100}{\giga\byte}}
    \label{tab:tpcds-size}
    \begin{tabularx}{\textwidth}{|X|c|c|c|}\hline
    \textbf{Формат} & \textbf{Объём, ГиБ} & \textbf{Число объектов} & \textbf{Отношение к ORC} \\ \hline
    ORC     & 22  & 618  & ${\times}1.00$ \\ \hline
    Parquet & 24  & 618  & ${\times}1.09$ \\ \hline
    Vortex  & 24  & 1206 & ${\times}1.09$ \\ \hline
    Lance   & 176 & 1302 & ${\times}8.00$ \\ \hline
    \end{tabularx}
\end{table}

Измерить пропускную способность записи на этом наборе не удалось: загрузка крупных
объёмов упиралась в пропускную способность объектного хранилища, и сопоставление
форматов по записи было бы некорректным. Этот аспект покрывается второй серией
экспериментов, где запись выполняется потоком с ограничением скорости.
